\section{Symbols, Indexing, and Sign Convention}

\subsection{Electrical Bases}
\begin{align}
S_{base} &= \text{system apparent-power base (MVA)}, \\
V_{base} &= \text{nominal line-line voltage base (kV)}, \\
Z_{base} &= \frac{V_{base}^2}{S_{base}}, \\
I_{base} &= \frac{S_{base}}{\sqrt{3}V_{base}}.
\end{align}

\subsection{Sets and Indexing}

\begin{align}
\mathcal{N}_{ac} &:\text{ AC buses}, &
\mathcal{E}_{ac} &:\text{ AC branches}, \\
\mathcal{N}_{dc} &:\text{ DC buses}, &
\mathcal{E}_{dc} &:\text{ DC branches}, \\
\mathcal{N}_{sl} &:\text{ AC slack buses}, &
\mathcal{N}_{PQ} &:\text{ AC PQ buses}, \\
\mathcal{N}_{PV} &:\text{ AC PV buses}, &
\mathcal{N}_{sl}^{\qual{dc}} &:\text{ DC voltage-controlled buses}.
\end{align}

The public model stores external bus identifiers in each component
(\texttt{bus}, \texttt{from\_bus}, \texttt{to\_bus}).  Most solver kernels then remap
those identifiers into contiguous zero-based local indices.

\subsection{Core AC and DC State Variables}
\begin{align}
V_i &= v_i e^{j\theta_i}, \\
Y_{ij} &= G_{ij} + jB_{ij}, \\
V_k^{\qual{dc}} &= v_k^{\qual{dc}}, \\
P_k^{\qual{dc,calc}} &= v_k^{\qual{dc}}(G_{dc}v^{\qual{dc}})_k.
\end{align}

For the main Newton power-flow path, the state vector is
\begin{align}
x_{PF} = \begin{bmatrix} \theta_{\mathcal{N}_{ac}\setminus\mathcal{N}_{sl}} & v_{\mathcal{N}_{PQ}} & v^{\qual{dc}}_{\mathcal{N}_{dc}\setminus\mathcal{N}_{sl}^{\qual{dc}}} \end{bmatrix}^{\top}.
\end{align}

For the native primal-dual AC OPF path, the decision vector is extended to
include generator dispatch and converter operating points:
\begin{align}
x_{OPF} = \begin{bmatrix}
\theta & v & P_g & Q_g & v^{\qual{dc}} & P^{\qual{conv,ac}} & Q^{\qual{conv,ac}} & P^{\qual{conv,dc}}
\end{bmatrix}^{\top},
\end{align}
with additional load-shedding variables in the parity formulation.

\subsection{Unified Sign Convention: Bus Injection Positive}

The current C++ module uses one consistent notation in the PF kernels and the
notebook adopts the same rule:
\begin{quote}
\textbf{Positive active or reactive power at a bus means injection into the network from the component attached to that bus.}
\end{quote}

Loads therefore store positive demand values in their data model, but they enter
power-balance equations with a minus sign.

\begin{longtable}{L{0.24\linewidth}L{0.48\linewidth}L{0.18\linewidth}}
\toprule
Quantity / field & Positive value means & Consequence in equations \\
\midrule
\endhead
\texttt{Generator::pg\_mw}, \texttt{qg\_mvar} & generator injects active/reactive power into its AC bus & added to \(P_i^{\qual{spec}}, Q_i^{\qual{spec}}\). \\
\texttt{Load::p\_mw}, \texttt{q\_mvar} & consumption requested from the network & subtracted from \(P_i^{\qual{spec}}, Q_i^{\qual{spec}}\); ZIP dependence scales the subtraction. \\
\texttt{Storage::p\_mw} & discharge to the grid; charging is negative \(P\) & treated as generation when positive, load when negative. \\
\texttt{Microgrid::p\_exchange\_mw} & export from the microgrid to the host AC grid & added as AC-bus injection in PF; ignored when the microgrid is islanded. \\
\texttt{VSCConverter} AC-side \(P^{\qual{conv,ac}}\) & converter injects active power into the AC bus & positive AC injection must be balanced by negative DC-bus injection plus losses. \\
\texttt{VSCConverter} DC-side \(P^{\qual{conv,dc}}\) & converter injects active power into the DC bus & the current code enforces \(P^{\qual{conv,ac}}+P^{\qual{conv,dc}}+P^{\qual{loss}}=0\). \\
\texttt{DCDCConverter::p\_ref\_mw} & positive transfer from \texttt{bus\_in} toward \texttt{bus\_out} & implemented as negative injection at \texttt{bus\_in} and positive injection at \texttt{bus\_out}. \\
\texttt{BranchFlow::pf\_mw} & active power sent from \texttt{from\_bus} to \texttt{to\_bus} & reporting only; no sign reversal in stored result. \\
\bottomrule
\end{longtable}

\subsection{Power-Balance Template}

With bus-injection-positive notation, the balance equations are written as
\begin{align}
\Delta P_i &= P_i^{\qual{spec}} - P_i^{\qual{calc}}, \\
\Delta Q_i &= Q_i^{\qual{spec}} - Q_i^{\qual{calc}}, \\
\Delta P_k^{\qual{dc}} &= P_k^{\qual{dc,spec}} - P_k^{\qual{dc,calc}}.
\end{align}

The specified term is assembled by adding all component injections and then
subtracting all demand terms.  This is the convention followed by the PF
residual evaluator, the native AC OPF assembly, and the parity formulation.
